% !TeX root = ../thuthesis-example.tex

% 中英文摘要和关键字

\begin{abstract}
针对中文方言语音处理中普遍存在的AI幻觉、数据稀缺及个性化学习体验缺失等挑战，本文提出并设计了一个面向多方言智能学习的“声学多模态认知校准框架”。该框架以构建的多模态方言音系知识图谱为认知锚点，通过融合声学、音系与语义特征并进行跨通道一致性评估，对声学模型的输出进行实时验证与校准，从而在有效抑制幻觉的同时，显著提升方言识别的准确率与鲁棒性。

在此核心框架之上，本项目进一步研发并集成了三个核心创新算法：“文脉智能体协作框架 (Dialectus)”基于多模态方言音系知识图谱，运用多智能体与思维链推理，实现了对复杂方言知识的深度问答；“自适应多步检索生成模型”通过动态建模学习者状态，动态查询重构，实现了高效的个性化教学；“音系自进化生成与校准模型”则构建了数据-反馈-优化的自进化闭环，以低成本方式完成了模型校准与自我迭代。

本项目将所有模块整合为一个功能完备的方言智能学习平台，旨在为方言的保护、传承与学习提供一个更准确、可靠、高效且智能化的解决方案。
  \thusetup{
    keywords = {方言音系知识图谱, 多智能体, 思维链, 自适应学习, 自进化},
  }
\end{abstract}


\begin{abstract*}
To address the prevalent challenges of AI hallucination, data scarcity, and the lack of personalized learning experiences in Chinese dialect speech processing, this paper proposes and designs a "Phono-Cognitive Alignment Framework" for intelligent multi-dialect learning. Grounded in a constructed multimodal phonological knowledge graph as a cognitive anchor, the framework performs real-time verification and calibration of the acoustic model's output by fusing acoustic, phonological, and semantic features and evaluating their cross-channel consistency. This approach not only effectively suppresses hallucinations but also significantly improves the accuracy and robustness of dialect recognition.

Building upon this core framework, this research further develops and integrates three core innovative algorithms: the "Dialectus" collaborative agent framework, which leverages the multimodal phonological knowledge graph, utilizes multi-agent systems and Chain-of-Thought reasoning to achieve in-depth question answering on complex dialectal knowledge; the "Adaptive Multi-step RAG model," which achieves efficient, personalized instruction by dynamically modeling the learner's state and performing dynamic query reconstruction; and the "Phonetic Self-Evolving Generation and Calibration Model," which establishes a self-evolving loop of data-feedback-optimization to complete model calibration and self-iteration in a low-cost manner.

This research integrates all modules into a fully functional intelligent dialect learning platform, aiming to provide a more accurate, reliable, efficient, and intelligent solution for the preservation, inheritance, and learning of dialects.
  \thusetup{
    keywords* = {dialect phonology knowledge graph, multi-agent, thought chain, adaptive learning, self-evolution},
  }
\end{abstract*}
